\chapter{Related Works}
\label{ch:related}

Related work falls into three parts: the LLM-based repair systems that
generate and validate candidate patches; the studies that measure what such
systems cost; and the techniques that make a repair loop cheaper. The techniques are grouped into five classes according to the part of the
cost they reduce. Test prioritisation, the group this thesis
works on, is reviewed in depth; the others briefly.

The review is a structured literature review, not a systematic one: a keyword
search of the major digital libraries and software-engineering venues for
LLM-based repair and for patch-validation cost, followed by backward and
forward snowballing from four seed
papers~\citep{qi2013frtp,lou2024aprregression,chen2021uniapr,xiao2024expressapr}.

\section{LLM-Based Repair Systems}
\label{sec:rel-systems}

A recent survey classifies LLM-based repair systems by how much control the
model has over the repair loop~\citep{yang2025llmrepairsurvey}; they also
differ in which part of the loop they try to improve.
Table~\ref{tab:rel-systems} groups representative systems into three families.

\begin{table}[htbp]
\centering
\caption{Representative LLM-based repair systems, what they optimise, and how
they use validation.}
\label{tab:rel-systems}
\footnotesize
\begin{tabularx}{\textwidth}{@{}>{\raggedright\arraybackslash}p{0.15\textwidth} >{\raggedright\arraybackslash}p{0.21\textwidth} >{\raggedright\arraybackslash}X >{\raggedright\arraybackslash}X@{}}
\toprule
Family & Systems & Optimises & Use of validation\\
\midrule
Conversational & ChatRepair~\citep{xia2024chatrepair} &
patch quality from failure feedback, at low cost per bug &
every candidate is tested; failures return to the conversation\\
Autonomous agents & RepairAgent~\citep{bouzenia2025repairagent},
SWE-agent~\citep{yang2024sweagent}, AutoCodeRover~\citep{zhang2024autocoderover} &
localisation and context gathering through tools &
the agent decides when to run tests\\
Agentless pipelines & Agentless~\citep{xia2025agentless} &
a fixed localise--repair--validate pipeline with low, predictable cost &
selected regression tests and generated reproduction tests filter and rank
candidates\\
\bottomrule
\end{tabularx}
\end{table}

\paragraph{Conversational and pipeline repair.} \textsc{ChatRepair}~\citep{xia2024chatrepair}
alternates generation and validation in one conversation: earlier incorrect
patches and their test failures stay in the prompt, so the model can avoid
repeating them. It reported fixing $162$ of $337$ bugs at about \$0.42 per bug,
yet $65\,\%$ of the patches it generated were duplicates of earlier
ones~\citep{chen2025conversational}. \textsc{Agentless}~\citep{xia2025agentless}
replaces free choice of actions with fixed phases (localisation, sampling of
several candidates and validation), which gives predictable cost and far fewer
tokens than tool-driven agents~\citep{yang2025llmrepairsurvey}.

\paragraph{Agent-, search- and reasoning-based repair.} Autonomous agents such as
\textsc{RepairAgent}~\citep{bouzenia2025repairagent},
\textsc{SWE-agent}~\citep{yang2024sweagent} and
\textsc{AutoCodeRover}~\citep{zhang2024autocoderover} let the model choose tools
to read, edit and test code; they improve localisation and context at the price
of long sessions, and \textsc{RepairAgent} spends $99\,\%$ of its repair time
executing tools. Cost figures appear for only a few
systems~\citep{yang2025llmrepairsurvey}.

\section{Validation in Current Tools}
\label{sec:rel-practice}

How current tools validate, deduplicate
and stop was read from the code of six open-source repair agents and from the
reports of deployed systems (Table~\ref{tab:rel-practice}). The review leads to four findings. No tool orders tests by how likely they are to fail or what they cost, and
none stops at the first failing test; validation is left to the model, run as
one fixed command, or run per patch after generation. Only
\textsc{AutoCodeRover} caches validation results, and only by exact text;
\textsc{Agentless} normalises patches, but only to vote among them. All tools set budgets in dollars, steps, rounds or timeouts, but only OpenHands
detects a stuck agent, from repeated actions and errors. No tool reports why
it gave up beyond an exit status, except an exit summary in Meta's repair
agent~\citep{maddila2025agentic}; \textsc{SWE-agent} submits its current diff
anyway. Every tool tracks tokens and money; none accounts for the time spent
running tests, although \textsc{RepairAgent} spends $99\,\%$ of its time in
tools~\citep{bouzenia2025repairagent}.

\begin{table}[htbp]
\centering
\caption[Validation, deduplication and stopping in current repair tools]{How
current repair tools validate, deduplicate and stop, read from their code
(open-source agents, November 2026) or reports (deployed systems).}
\label{tab:rel-practice}
\scriptsize
\begin{tabularx}{\textwidth}{@{}>{\raggedright\arraybackslash}p{0.15\textwidth} >{\raggedright\arraybackslash}X >{\raggedright\arraybackslash}p{0.17\textwidth} >{\raggedright\arraybackslash}p{0.24\textwidth}@{}}
\toprule
Tool & Validation & Repeats & Budget, stop, give-up\\
\midrule
\textsc{SWE-agent}~\citep{yang2024sweagent} & the model runs a reproduction script through the shell & identical parallel actions dropped & \$ and call limits, timeouts; on exit submits the current diff\\
OpenHands~\citep{openhands2026sdk} & left to the model & notice after repeated identical calls & step and \$ limits; stuck detector; status only\\
\textsc{Agentless}~\citep{xia2025agentless} & regression and reproduction tests on every sample, after generation; no feedback & normalised, for voting only & number of samples, timeouts\\
\textsc{AutoCode\-Rover}~\citep{zhang2024autocoderover} & reproducer, an LLM reviewer, then the developer's tests & exact-text cache of results & retry and round caps; a reason string\\
Aider~\citep{aider2026} & linter, then the user's test command & notice for an edit already applied & three reflections\\
\textsc{RepairAgent}~\citep{bouzenia2025repairagent} & the whole Defects4J suite & notice for a repeated command & 40 cycles\\
\midrule
SapFix (Meta)~\citep{marginean2019sapfix} & reproduction, analysis, CI tests & one fix published & dropped if unreviewed in 7 days\\
Repair agent (Meta)~\citep{maddila2025agentic} & the broken tests, analysis, an LLM judge & failures grouped by cause & timeout and cap; exit summary\\
Passerine (Google)~\citep{rondon2025passerine} & the bug-reproduction tests & --- & 20 trajectories of 25 steps\\
\bottomrule
\end{tabularx}
\end{table}

\section{Cost Measurement in LLM-Based Repair}
\label{sec:rel-cost}

Studies that report what LLM-based repair costs fall into three kinds.

\paragraph{Cost per bug, as a by-product.} Repair systems report a cost next to
their success rate, each in its own unit: \textsc{ChatRepair} in API dollars per
bug~\citep{xia2024chatrepair}, \textsc{RepairAgent} in tokens and in repair time
per bug~\citep{bouzenia2025repairagent}, and \textsc{CigaR}, which makes token
cost its objective, in tokens per bug against a conversational
baseline~\citep{hidvegi2024cigar}. These figures compare systems on one side of
the bill and are not normalised across studies~\citep{yang2025llmrepairsurvey}.

\paragraph{Cost-aware comparisons.} A second kind of study holds the budget
fixed and asks which strategy spends it better. Independent repeated prompting repairs
more bugs than a conversational repairer at the same number of
attempts~\citep{chen2025conversational}. For software agents,
\textsc{SWE-Effi} re-scores SWE-bench systems by their resolve rate per token
and per unit of time, and finds that effectiveness depends on how a scaffold
pairs with its model, and that unresolved issues consume far more resources
than resolved ones~\citep{fan2025sweeffi}.

\paragraph{The cost of execution.} \citet{lin2026torun} isolate execution in
agents. Every agent examined relies on code execution, which has become
standard practice, but applies it indiscriminately: for commercial agents on
SWE-bench, prohibiting execution changed the resolve rate by only $1.25$
percentage points, a difference that was not statistically significant, while
saving $56$--$62\,\%$ of tokens and $48$--$54\,\%$ of wall-clock time for one of
them.

\paragraph{What is not measured.} All three kinds treat a repair session as a
whole, or remove one of its parts entirely. None splits the cost of one loop
into generation and validation, nor validation into repeated, rejected and
accepted candidates (Section~\ref{sec:bg-cost}). Without that split, the share of
a session that a validation-cost technique can reach is unknown, and so is the
largest saving it could deliver.

\section{Classification of Cost-Reduction Techniques}
\label{sec:rel-taxonomy}

Table~\ref{tab:rel-taxonomy} sorts the techniques into five groups by the part
of Equation~\eqref{eq:bg-decomp} they cut. Because the groups cut different
parts, they can be combined. Groups A and B are the two classical problems of
regression testing: \emph{selection} decides which tests to run,
\emph{prioritisation} the order in which they run.

\begin{table}[htbp]
\centering
\caption[The five groups of cost reduction]{The five groups of cost reduction.
``Learns'' marks groups that use what the session has observed.}
\label{tab:rel-taxonomy}
\footnotesize
\setlength{\tabcolsep}{4pt}
\begin{tabularx}{\textwidth}{@{}>{\raggedright\arraybackslash}p{0.18\textwidth} >{\raggedright\arraybackslash}p{0.19\textwidth} >{\raggedright\arraybackslash}p{0.08\textwidth} >{\raggedright\arraybackslash}p{0.13\textwidth} c >{\raggedright\arraybackslash}X@{}}
\toprule
Group & Cuts & Safe & Language-bound & Learns & Representative work\\
\midrule
A test selection & tests that do not reach the change & if sound & yes & no & RTS in APR~\citep{lou2024aprregression}\\
\textbf{B test prioritisation} & $V_{\mathrm{rej}} - V_{\mathrm{rej}}^{*}$ & yes & no & yes & \frtp~\citep{qi2013frtp}, PIT~\citep{coles2016pit}\\
C execution acceleration & start-up share & yes & yes & no & UniAPR~\citep{chen2021uniapr}, ExpressAPR~\citep{xiao2024expressapr}\\
D patch reduction & $V_{\mathrm{rep}}$ and equivalent candidates & mostly & partly & partly & verdict cache, test-equivalence~\citep{mechtaev2018testequivalence}\\
E generation-side reduction & $G_{\mathrm{rep}}$ and rounds & --- & no & yes & ChatRepair~\citep{xia2024chatrepair}, CigaR~\citep{hidvegi2024cigar}\\
\bottomrule
\end{tabularx}
\end{table}

\section{Test Prioritisation}
\label{sec:rel-b}

\paragraph{Principle.} Validation can stop as soon as one test fails, because
one failure proves the patch wrong. The order of the tests then decides how
much is spent before that failure: a failing test placed first costs one run,
the same test placed last costs the whole suite. Ordering removes no test, so
for deterministic tests it never changes a verdict.

\paragraph{Key techniques.} The information an order is built from
(Table~\ref{tab:rel-b}).

\begin{table}[htbp]
\centering
\caption{Test prioritisation, by the information it ranks on.}
\label{tab:rel-b}
\footnotesize
\begin{tabularx}{\textwidth}{@{}>{\raggedright\arraybackslash}p{0.17\textwidth} X >{\raggedright\arraybackslash}p{0.27\textwidth}@{}}
\toprule
Information & Rule & Representative work\\
\midrule
Coverage & tests covering more (new) code first & \tptot, \tpadd~\citep{rothermel2001prioritizing}\\
Failure history & tests that rejected more candidates first & \frtp~\citep{qi2013frtp}\\
Cost & likelihood of failing per unit of cost & least-cost testing~\citep{mitten1960}, APFD$_C$~\citep{elbaum2001cost}\\
Location tiers & failing tests, then tests in the package of the patched code, then the rest & ExpressAPR~\citep{xiao2024expressapr}\\
Coverage, cost and name & tests covering the change only, fastest first, the unit test of the changed class favoured & PIT~\citep{coles2016pit}\\
History across variants & tests that killed other mutants first & FaMT~\citep{zhang2013famt}\\
History across CI cycles & a learned policy, rewarded for early failures & RETECS~\citep{spieker2017retecs}, COLEMAN~\citep{lima2022coleman}\\
\bottomrule
\end{tabularx}
\end{table}

Coverage-based
ordering comes from regression testing (Section~\ref{sec:bg-regression}) and
reads only the program. \frtp\ is built for repair: it runs first the tests
that rejected earlier candidates in the same session, information that validation provides at no extra cost, and it matched GenProg's success rate while
validating faster~\citep{qi2013frtp}. Cost entered prioritisation in regression
testing~\citep{elbaum2001cost}, where test costs are assumed known in advance;
none of the repair orderings uses it, although the run of the faulty program
before repair already measures each test's cost. ExpressAPR orders tests by location in fixed tiers: the failing tests, then
the tests in the package of the patched code, then the
rest~\citep{xiao2024expressapr}. It uses neither test costs nor the session's
history.

Mutation testing faces the same problem as repair: many variants of one
program, each to be rejected by any failing test, a duality ExpressAPR builds
on~\citep{xiao2024expressapr}. PIT, the most used mutation tool for Java, runs
only the tests that cover the mutated line, fastest first, favours the unit
test of the mutated class, and stops at the first kill~\citep{coles2016pit}.
FaMT orders the tests for each mutant by how they fared on other
mutants~\citep{zhang2013famt}. Together they use the signals a repair ordering could use, namely cost, the
location of the change and history across variants, but in a setting where the variants are known in advance, the
change is located exactly by an instrumented coverage run, and no failure is
reported to anyone. In LLM-based repair the candidates arrive one at a time and
are few, and the failing test that validation reports becomes the model's next
input.

Learned orders exist in continuous integration, where the order is learned
across commits: RETECS uses reinforcement learning~\citep{spieker2017retecs}
and COLEMAN a multi-armed bandit~\citep{lima2022coleman}, both rewarded for
detecting failures early. Ranking under stop-at-first-success has its own bandit
model, cascading bandits, with UCB and Thompson sampling
algorithms~\citep{kveton2015cascading,zhong2021tscascade} and a variant in
which examining each item has a cost~\citep{zhou2018costcascading}. These
methods learn over hundreds or thousands of cycles. A repair session validates
a handful of candidates, so whatever a learned order explores must pay off
within a few rounds.

\paragraph{Limitations.} \frtp\ ranks on a \emph{count} of past
rejections~\citep{qi2013frtp}, which cannot tell a test never run from one run
often that never rejected, and no repair ordering reads how long a test takes;
mutation tools do, but statically, from an instrumented run, and with rules
that join the signals in sequence rather than in one estimate.
History-based orders also learn from censored evidence: stopping at the first
failure records only one failing test per rejected patch
(Section~\ref{sec:bg-censoring}). Ordering saves anything only when validation
stops at the first failure, and after test selection there are fewer tests to
order, so the best order can change~\citep{lou2024aprregression}. Finally, the
unit of measurement matters: prioritisation studies count test executions,
which reproduce across machines~\citep{lou2024aprregression}, but when test
costs differ an order that runs more cheap tests can finish sooner than one
that runs fewer expensive ones, so counts and seconds can rank two orders in
opposite directions.

\section{Other Cost-Reduction Techniques}
\label{sec:rel-other}

\paragraph{A: Test selection.} A test that never executes the code a patch
modifies gives the same result on the patched program as on the original, so it
can be skipped. \citet{lou2024aprregression} measured regression test selection
inside twelve repair systems on more than two million patches: every granularity (class~\citep{gligoric2015ekstazi}, method or statement)
substantially improved efficiency, finer granularity was better, and combining
selection with prioritisation was best. Sound
selection needs a dependency tracker for the language, and it has nothing to
remove when every test runs the whole program.

\paragraph{C: Execution acceleration.} Part of a test run is overhead:
recompiling the patched program and starting the runtime. UniAPR keeps one JVM
alive and reloads only the patched classes, speeding repair up by over an order
of magnitude~\citep{chen2021uniapr}; ExpressAPR borrows five mutation-testing
accelerations, including compiling many candidates into one
program, and validates $137\times$ faster than plain
validation~\citep{xiao2024expressapr}. Both are built for the JVM; both also
prioritise tests, ExpressAPR with the static location tiers of
Table~\ref{tab:rel-b}.

\paragraph{D: Patch reduction.} Some candidates need not be run at all. The
simplest reduction is a verdict cache: a candidate textually identical to one
already validated is given the stored verdict, which matters for LLMs, since
$65\,\%$ of the patches \textsc{ChatRepair} generated were
duplicates~\citep{chen2025conversational}. Test-equivalence analysis runs one
candidate per group of candidates that behave alike on the
tests~\citep{mechtaev2018testequivalence}. Equivalence analysis needs a program analysis
for the language.

\paragraph{E: Generation-side reduction.} A loop that needs fewer rounds,
or wastes fewer on candidates it has already seen, is cheaper on both sides.
The main levers are what the model is shown after a failure, such as the
failing test in a conversation~\citep{xia2024chatrepair}, and how the budget is
spent: many
candidates per call to amortise the prompt~\citep{hidvegi2024cigar}, or
independent re-sampling instead of conversational
refinement~\citep{chen2025conversational}. These
techniques are evaluated by success and tokens; the validation they cause or
save is not reported.

\section{Research Gap}
\label{sec:rel-gap}

Three observations close the review. First, every validation-side group has a technique with a measured saving:
test selection~\citep{lou2024aprregression},
test prioritisation~\citep{qi2013frtp}, execution
acceleration~\citep{xiao2024expressapr} and equivalence-based patch
reduction~\citep{mechtaev2018testequivalence}. However, all were measured with
pre-LLM repair systems, on Java or C programs.

Second, none of the LLM-based systems of Section~\ref{sec:rel-systems} reports
using test prioritisation, and where these systems report cost, they report one
side of the bill in a unit of their own (Section~\ref{sec:rel-cost}). How the
cost of an LLM-based loop splits between generation and validation, and within
validation, is not known.

Third, of the validation-side groups, test prioritisation transfers most
directly (Table~\ref{tab:rel-taxonomy}): it is safe for deterministic tests,
language-independent, and learns from the loop, while execution acceleration
is bound to one runtime, test selection to a sound dependency tracker, and
patch reduction beyond a verdict cache to program analysis. But it acts only on
the cost of rejecting invalid patches, so its value depends on how large that cost is in the whole session, a
quantity that no study has measured.


\paragraph{Open problems.} Together, the observations give five open problems,
which the proposed system addresses (Chapter~\ref{ch:system}).
\begin{enumerate}[label=\textbf{O\arabic*},leftmargin=*,itemsep=2pt]
\item \textbf{The cost composition of a repair loop is unmeasured.} Without it,
the upper bound on what a technique can save in a given setting is unknown.
\item \textbf{The repair orderings were built for another kind of loop.} \frtp\
was designed with search-based generators, which validate many
candidates~\citep{qi2013frtp}; an LLM-based loop validates
few, so a session-local history is short, and in a conversational loop the
failing test is also the model's next input.
\item \textbf{Count instead of rate, under censored evidence.} A count of past
rejections cannot tell a test never run from one that never rejects, and early
stopping hides every test after the first failure.
\item \textbf{Cost and the location of the patch are used only statically.}
The loop times every test it runs and knows which class each candidate
modifies; no repair ordering uses the costs, and the location enters repair
orderings only as a fixed tier (ExpressAPR). Mutation tools combine cost and
location, but statically and from an instrumented run (PIT); none estimates
them online from the session's own runs.
\item \textbf{Repeats and stuck sessions are not managed.} Most of the patches
a conversational repairer generated were repeats of earlier
ones~\citep{chen2025conversational}; they are generated, and often validated,
again, since only one tool caches results, by exact text. Sessions stop at a
fixed budget; only one agent framework detects that it is stuck, and none
reports why it gave up (Table~\ref{tab:rel-practice}).
\end{enumerate}
